% Chapter 5 -- Building programs (model: TP 6.0 UG Ch 6 "Project
% management": make and build become the compile-assemble-link cycle of a
% single-file program).  AN-19.
% Sources (Law 1): os/tps/tps.pas RunCodeGenerator (TPSOUT.S rewritten only
% by a clean fourth pass); the Makefile recipes of test-tps-gen-fpc and
% test-tps-gen-os (nasm, ld with seed/rt/seed_dos.ld + start_dos + fileio
% objects, objcopy to a flat .COM; TPS.COM run on InitechDOS, TINY.COM run
% on InitechDOS); spec/memory_map.h (image area 40100h..6F000h, 64 KiB
% stack); os/tps/CODEGEN.md s9 (the compiler's own size).
\Gates{test-tps-gen-fpc, test-tps-gen-build, test-tps-gen-os, test-compiler-os, test-pubs-turbo}
\chapter{Building programs}
\ix{building programs}

This chapter describes the steps between a program that compiles and a
program file you can run, and the limits a program must keep within.

\section{Compiling under InitechDOS}
\ix{TPS.COM}
To compile the program in \file{PAYROLL.PAS}:

\begin{enumerate}
\item Type \typed{COPY PAYROLL.PAS TPSIN.PAS} and press \key{Enter}.
\item Type \typed{TPS} and press \key{Enter}.
\item Check that the output on COM1 ends with \code{TPS-GEN-OK}.
\end{enumerate}

You can put the two commands into a batch file (\doctitle{InitechDOS
User's Guide}, Chapter 5) so that one command compiles any program.

\ArrowNote{The compiler writes \file{TPSOUT.S} only when the whole program
has compiled.  After an error the \file{TPSOUT.S} of an earlier
compilation is left as it was; always check for \code{TPS-GEN-OK} before
you use it.}

\section{From TPSOUT.S to a program file}
\ix{TPSOUT.S}
\ix{.COM file}
\file{TPSOUT.S} is assembled into an object module, and the object module
is linked with the objects of the Turbo Initech run-time library: the
InitechDOS start object, which must come first, and, when the program uses
files, the file object.  The
linked program, placed at address 40100h, is written out as a flat binary
image with the extension \file{.COM}.  Give it the name you want to type at
the prompt: \file{PAYROLL.COM} is run by typing \typed{PAYROLL}.
\draft{the assembler and linker run on the development system (Self-
Compilation Guide, Chapter 3); an InitechOS-resident assembler and linker
are milestone M8, bead initech-ls4}

\section{Limits of a program}
\ix{limits}
\ix{memory, limits}
\begin{BorTable}{Program limits}{p{36mm}p{62mm}}
\Hd{Item} & \Hd{Limit}\EndHd
Code, constants and global variables together & from 40100h up to 6F000h,
  about 187K\\
Stack: frames of active routines and temporary values & 64K\\
Source file & 262,144 characters\\
One string & 255 characters\\
\end{BorTable}

Global variables count toward the first limit; a program with large arrays
uses its space quickly, since every element of an array of integers takes
four bytes.  Deep recursion uses the stack: each active call takes the
size of its frame, its arguments and its temporary values.  The compiler's
own limits on the number of names and symbols are in Chapter 8.

\section{Building Turbo Initech itself}
The compiler is a Turbo Initech program: its source is one Pascal file,
and it is compiled exactly as this chapter describes.  The
\doctitle{Self-Compilation Guide} describes how the first compiler is
built and how a compiler proves that it reproduces itself.
